% Corrections and Concessions
% Recovered from three-possibilities-interlude.tex (pre-Phase-2b)
% Plan 0094: relocated to appendix as standalone reference

\chapter*{Corrections and Concessions}
\addcontentsline{toc}{chapter}{Corrections and Concessions}

The source material in this book spans twenty years. Some of my earlier writing contains errors --- factual assertions I stated with confidence that turned out to be wrong, or details that drifted in my memory over two decades. A careful reader, or a hostile one, will find them.

This section lists the errors I have identified as of February 2026. I prefer you find them here rather than discover them yourself and wonder what else I got wrong.

\textbf{1. Bravo Two Zero patrol composition.} In my earlier documents, I stated the B20 patrol had six members and that two escaped (Chris Ryan and Healer). All published accounts --- McNab, Ryan, Asher, Coburn --- agree the patrol had eight members and that one escaped (Ryan). No published source supports either a six-man patrol or a second escapee. SAS does not publish official operational histories, so the published accounts are contested personal narratives, not authoritative records. But the weight of available evidence contradicts my assertion. I acknowledge this error. The core proposition --- that Healer was in the patrol --- is separate from the patrol size and is not affected by this correction.

\textbf{2. The Coventry myth.} In earlier drafts, I repeated F.W.\ Winterbotham's claim that Churchill had specific advance intelligence that Coventry was the Luftwaffe target on November 14, 1940, and chose not to act in order to protect Ultra. The majority of modern historians (Hinsley, R.V.\ Jones, GCHQ's own published account) reject this. Churchill may have known a major raid was coming but probably did not know Coventry was the specific target. The corrected version appears in the chapter on the Code War. The broader point --- that Ultra secrecy was considered so important that historians take seriously whether a leader would sacrifice a city to protect it --- stands. The specific anecdote I used to illustrate it was wrong.

\textbf{3. Obscure Images identification.} I previously identified Healer's Cult of the Dead Cow handle as ``Obscure Images.'' Stylometric analysis of the Obscure Images corpus revealed an art-saturated voice from Northern Illinois University --- a consistent creative identity spanning 1989--1996. When confronted with this, I confirmed that Healer was ``totally not art-aware.'' The dominant Obscure Images voice is not Healer. The handle may have been shared (cDc is known for shared handles --- Oxblood Ruffin is publicly admitted as one), or I may have been wrong about the identification entirely.

\textbf{4. Churchill / Coventry as analogy precedent.} See \#2. I used this as a key analogy for ``what governments do to protect cryptographic secrets.'' The analogy is valid --- the Coventry \textit{debate} proves the point --- but I stated the myth as near-fact in early drafts. Corrected.

\textbf{5. King Fahd / King Faisal.} In verbal accounts of the Battle of Baghdad, I said ``King Faisal'' called President Bush to cancel the Iraq invasion. King Faisal died in 1975. The correct name is King Fahd bin Abdulaziz Al Saud. My archive has the correct name. This was a memory slip in conversation, not an error in my written documents.

\textbf{6. ``David Lane'' as Healer's name.} My earlier documents use the name ``David Lane'' for Healer. Lane is not his real name. I used it in my documents for the same reason I use ``Healer'' here --- to tell the story without exposing his actual identity. Where earlier documents present ``Lane'' as a real name, that was a pseudonym.

\textbf{7. EO 13026 timeline.} In earlier analysis, I implied the November 1996 executive order relaxing PKC export controls was a direct response to Shor's 1994 algorithm. The classified program predated Shor by approximately five years. The correct parallel is structural: GCHQ had PKC \textasciitilde{}1973, public discovery 1976--78 (gap: \textasciitilde{}4 years). Classified program had quantum factoring \textasciitilde{}1987--1989, Shor published 1994 (gap: \textasciitilde{}5--7 years). Shor's role maps to Diffie-Hellman's: independent public rediscovery of something that already existed in classified channels. The EO 13026 connection needs re-verification.

\vspace{1cm}

\textbf{What these errors tell you.} Under Option A (confabulation), these errors are cracks in a fabrication --- details the fabricator didn't research carefully enough. Under Option B (exaggerated kernel), they are the kind of drift you'd expect as a real story grows and blurs over twenty years. Under Option C (substantially true), they are memory degradation: the architecture of events stays intact while specific details shift. All three possibilities predict errors. The question is whether the errors are structural (contradicting the core narrative) or peripheral (details that drifted). Every error listed above is peripheral. None contradicts the core propositions: the program, the technology, the relinquishment, the Guardian.

If I had perfect recall of every detail from twenty years ago, that would be more suspicious, not less.
